% =====================================================================
%  macros.tex
%  Mathematical shorthands. Keep this short; add as the book grows.
% =====================================================================

% --- Number sets ------------------------------------------------------
\newcommand{\R}{\mathbb{R}}
\newcommand{\Z}{\mathbb{Z}}
\newcommand{\N}{\mathbb{N}}
\newcommand{\Q}{\mathbb{Q}}
\newcommand{\Cplx}{\mathbb{C}}   % \C is sometimes claimed by other packages

% --- Vectors ----------------------------------------------------------
% Default style: bold roman. Switch to \boldsymbol if you prefer italic.
\newcommand{\vect}[1]{\mathbf{#1}}
\newcommand{\zerovec}{\mathbf{0}}

% --- Partial derivatives ---------------------------------------------
\newcommand{\pd}[2]{\dfrac{\partial #1}{\partial #2}}
\newcommand{\pdmix}[3]{\dfrac{\partial^{2} #1}{\partial #2 \, \partial #3}}

% --- Norms and absolute value ----------------------------------------
\DeclarePairedDelimiter{\norm}{\lVert}{\rVert}
\DeclarePairedDelimiter{\abs}{\lvert}{\rvert}

% --- Differential forms ----------------------------------------------
% Use \dform for the upright "d" if you want forms typeset distinctively.
% (Many authors leave it italic; switch the definition to taste.)
\newcommand{\dform}{\mathrm{d}}

% --- Common operators -------------------------------------------------
\DeclareMathOperator{\grad}{grad}
\DeclareMathOperator{\divg}{div}
\DeclareMathOperator{\curl}{curl}
\DeclareMathOperator{\Hess}{Hess}
\DeclareMathOperator{\Image}{Im}
\DeclareMathOperator{\Kernel}{Ker}
\DeclareMathOperator{\spn}{span}
\DeclareMathOperator{\sgn}{sgn}

% --- Inner products and pairings -------------------------------------
\newcommand{\inner}[2]{\langle #1, #2 \rangle}
